% p_n ; (id^{\otimes(n-1)} \otimes !) = p_{n-1}
% Read top to bottom. Discarding the n-th wire of p_n yields p_{n-1}.
\documentclass[tikz,border=0pt]{standalone}
\usepackage{fritzstrings}
\begin{document}
\begin{tikzpicture}[x=1cm,y=1cm]
  \begin{scope}[xshift=-3.15cm]
    \draw[fill=white] (-0.7,3.15) -- (0.7,3.15) -- (1.85,1.85) -- (-1.85,1.85) -- cycle;
    \node at (0,2.28) {$p_n$};
    \draw (-0.85,1.85) -- (-0.85,-1.35);
    \node[bn] (d) at (0.85,-0.15) {};
    \draw (0.85,1.85) -- (d);
    \node[wl,right] at (0.95,0.85) {$X$};
    \node[wl,right] at (1.02,-0.15) {$!$};
    \node[wl,below] at (-0.85,-1.35) {$X^{\otimes(n-1)}$};
    \node[wl,below] at (-0.85,-2.05) {$\mathrm{id}^{\otimes(n-1)}$};
  \end{scope}

  \node[eq] at (0,0.55) {$=$};

  \begin{scope}[xshift=2.85cm]
    \draw[fill=white] (-0.7,3.15) -- (0.7,3.15) -- (1.35,1.85) -- (-1.35,1.85) -- cycle;
    \node at (0,2.28) {$p_{n-1}$};
    \draw (0,1.85) -- (0,-1.35);
    \node[wl,below] at (0,-1.35) {$X^{\otimes(n-1)}$};
  \end{scope}

  \node[font=\small] at (0,-3.15)
    {$p_n\,;\,\bigl(\mathrm{id}^{\otimes(n-1)}\otimes !\bigr)=p_{n-1}$};
\end{tikzpicture}
\end{document}
